% GENERATED FILE. Edit canonical lessons, then run npm run generate:books.
% canonical-source-hash: fnv1a64:fa8a10fa0af4a8c4
% canonical-lessons: GE-C251-uebrigens, GE-C251-sowieso, GE-C251-hoechstens, GE-C251-mindestens, GE-C251-waehrend, GE-R251-first-pass-words-from-chapters-247-248, GE-R251-second-pass-words-from-chapters-249-251

\chapter{By the way, Anyway, At most, At least, During}
\label{ch:ge-by-the-way-anyway-at-most-at-least-during}
\hlchaptermodality{251}

\begin{chapteropening}
\textbf{By the end of this chapter:} \emph{I can say by the way, anyway, at most, at least and during.}

Complete the last lesson of chapter 251: i can say by the way, anyway, at most, at least and during.
\end{chapteropening}

\section[übrigens]{übrigens — by the way}

\label{lesson:GE-C251-uebrigens}



\begin{quote}
Before the new one: say the German for in any case, then the German for actually.
\end{quote}

\subsection*{You'll want to know: übrigens}
\textbf{übrigens} — \textquotedblleft{}by the way\textquotedblright{}.

\begin{grammarlens}[title={What you've built}]
One more word for this chapter.
\end{grammarlens}

\subsection*{Your turn}
\begin{itemize}
\raggedright
  \item \emph{Say it:} \emph{übrigens}
  \item \emph{Say it:} \emph{übrigens}, once more
  \item \emph{Say it:} say \emph{übrigens}
  \item \emph{Recall:} say the German for in any case, then the German for actually, then say \emph{übrigens} again
\end{itemize}

\begin{tcolorbox}[breakable,colback=teal!4,colframe=teal!35!black,title={Before you move on}]
What does übrigens mean? (\textquotedblleft{}By the way\textquotedblright{}.) Say it once more.
\end{tcolorbox}
\section[sowieso]{sowieso — anyway}

\label{lesson:GE-C251-sowieso}



\begin{quote}
Before the new one: say the German for actually, then the German for by the way.
\end{quote}

\subsection*{You'll want to know: sowieso}
\textbf{sowieso} — \textquotedblleft{}anyway\textquotedblright{}.

\begin{grammarlens}[title={What you've built}]
One more word for this chapter.
\end{grammarlens}

\subsection*{Your turn}
\begin{itemize}
\raggedright
  \item \emph{Say it:} \emph{sowieso}
  \item \emph{Say it:} \emph{sowieso}, once more
  \item \emph{Say it:} say \emph{sowieso}
  \item \emph{Recall:} say the German for actually, then the German for by the way, then say \emph{sowieso} again
\end{itemize}

\begin{tcolorbox}[breakable,colback=teal!4,colframe=teal!35!black,title={Before you move on}]
What does sowieso mean? (\textquotedblleft{}Anyway\textquotedblright{}.) Say it once more.
\end{tcolorbox}
\section[höchstens]{höchstens — at most}

\label{lesson:GE-C251-hoechstens}



\begin{quote}
Before the new one: say the German for by the way, then the German for anyway.
\end{quote}

\subsection*{You'll want to know: höchstens}
\textbf{höchstens} — \textquotedblleft{}at most\textquotedblright{}.

\begin{grammarlens}[title={What you've built}]
One more word for this chapter.
\end{grammarlens}

\subsection*{Your turn}
\begin{itemize}
\raggedright
  \item \emph{Say it:} \emph{höchstens}
  \item \emph{Say it:} \emph{höchstens}, once more
  \item \emph{Say it:} say \emph{höchstens}
  \item \emph{Recall:} say the German for by the way, then the German for anyway, then say \emph{höchstens} again
\end{itemize}

\begin{tcolorbox}[breakable,colback=teal!4,colframe=teal!35!black,title={Before you move on}]
What does höchstens mean? (\textquotedblleft{}At most\textquotedblright{}.) Say it once more.
\end{tcolorbox}
\section[mindestens]{mindestens — at least}

\label{lesson:GE-C251-mindestens}



\begin{quote}
Before the new one: say the German for anyway, then the German for at most.
\end{quote}

\subsection*{You'll want to know: mindestens}
\textbf{mindestens} — \textquotedblleft{}at least\textquotedblright{}.

\begin{grammarlens}[title={What you've built}]
One more word for this chapter.
\end{grammarlens}

\subsection*{Your turn}
\begin{itemize}
\raggedright
  \item \emph{Say it:} \emph{mindestens}
  \item \emph{Say it:} \emph{mindestens}, once more
  \item \emph{Say it:} say \emph{mindestens}
  \item \emph{Recall:} say the German for anyway, then the German for at most, then say \emph{mindestens} again
\end{itemize}

\begin{tcolorbox}[breakable,colback=teal!4,colframe=teal!35!black,title={Before you move on}]
What does mindestens mean? (\textquotedblleft{}At least\textquotedblright{}.) Say it once more.
\end{tcolorbox}
\section[während]{während — during, while}

\label{lesson:GE-C251-waehrend}



\begin{quote}
Before the new one: say the German for at most, then the German for at least.
\end{quote}

\subsection*{You'll want to know: während}
\textbf{während} — \textquotedblleft{}during, while\textquotedblright{}.

\begin{grammarlens}[title={What you've built}]
One more word for this chapter.
\end{grammarlens}

\subsection*{Your turn}
\begin{itemize}
\raggedright
  \item \emph{Say it:} \emph{während}
  \item \emph{Say it:} \emph{während}, once more
  \item \emph{Say it:} say \emph{während}
  \item \emph{Recall:} say the German for at most, then the German for at least, then say \emph{während} again
\end{itemize}

\begin{tcolorbox}[breakable,colback=teal!4,colframe=teal!35!black,title={Before you move on}]
What does während mean? (\textquotedblleft{}During, while\textquotedblright{}.) Say it once more.
\end{tcolorbox}
\section[(dialogue)]{Words from chapters 247-248 — a first pass}

\label{lesson:GE-R251-first-pass-words-from-chapters-247-248}



\begin{quote}
Say the words for at least and during.
\end{quote}

\subsection*{Your turn}
Say these aloud:

\begin{itemize}
\raggedright
  \item to pass, to prove, to bloom, to discuss, to pack
  \item to furnish, to end, to discover, to excuse, to find out
\end{itemize}

\begin{tcolorbox}[breakable,colback=teal!4,colframe=teal!35!black,title={Before you move on}]
To pass? (\textbf{bestehen}.) To prove? (\textbf{beweisen}.) To bloom? (\textbf{blühen}.) To discuss? (\textbf{diskutieren}.) To pack? (\textbf{einpacken}.) To furnish? (\textbf{einrichten}.) To end? (\textbf{enden}.) To discover? (\textbf{entdecken}.) To excuse? (\textbf{entschuldigen}.) To find out? (\textbf{erfahren}.)
\end{tcolorbox}
\section[(dialogue)]{Words from chapters 249-251 — a second pass}

\label{lesson:GE-R251-second-pass-words-from-chapters-249-251}



\begin{quote}
Say the words for to complete and to get done.
\end{quote}

\subsection*{Your turn}
Say these aloud:

\begin{itemize}
\raggedright
  \item to complete, to get done, to appear, to produce, to demand
  \item besides, therefore, nevertheless, in any case, actually
  \item by the way, anyway, at most, at least, during
\end{itemize}

\begin{tcolorbox}[breakable,colback=teal!4,colframe=teal!35!black,title={Before you move on}]
To complete? (\textbf{ergänzen}.) To get done? (\textbf{erledigen}.) To appear? (\textbf{erscheinen}.) To produce? (\textbf{erzeugen}.) To demand? (\textbf{fordern}.) Besides? (\textbf{außerdem}.) Therefore? (\textbf{deshalb}.) Nevertheless? (\textbf{trotzdem}.) In any case? (\textbf{jedenfalls}.) Actually? (\textbf{eigentlich}.) By the way? (\textbf{übrigens}.) Anyway? (\textbf{sowieso}.) At most? (\textbf{höchstens}.) At least? (\textbf{mindestens}.) During? (\textbf{während}.)
\end{tcolorbox}
